\frontchapter{摘要}
近似最近邻（ANN）搜索支持大规模嵌入向量集合上的相似性检索。GPU 可以加速搜索，但有限显存使压缩索引具有吸引力。JHQ 是一种层级量化方法：先用紧凑编码筛选候选，再利用残差编码对较小的候选集精排。本文开发 JHQ-GPU，研究如何利用这种表示实现高效的 GPU 建库与搜索。

实现将 GPU 模型拟合和分批编码，与有界输入暂存、打包且按倒排列表排列的布局结合。扫描将每张 256 项距离表替换为两张 16 项表，保持实数算术下的分数。分组扫描在固定的已计算分数与位置排序下，保持精确的全局候选选择。校准通过比较当前结果与较大预算参考返回的近邻，在无需真值标签的情况下决定精排候选数量。

在 NVIDIA RTX~5090 上，跨方法评测使用 Vogue、arXiv、OpenAI-1536 和 OpenAI-3072。在召回率 0.97、batch 1,000 下，JHQ-GPU 在这四个数据集上的主机到主机吞吐为图索引 CAGRA fp32 的 $1.32$--$1.75\times$，并在其中三个数据集上取得已测方法的最高吞吐。对于更大规模的 BGE 和 Stella 文本嵌入数据集，原始 FP32 向量的存储需求已超出设备显存，JHQ-GPU 仍可在单张 GPU 上完成搜索；最大索引包含 17.78 百万个向量，常驻 GPU 显存为 20.44\,GiB。
